\subsection{Productos tensoriales completados}
\label{subsection:0.7.7-fr}

\begin{env}[7.7.1]
\label{0.7.7.1-fr}
Sean $A$ un anillo linealmente topologizado, $M$, $N$, dos $A$-módulos
linealmente topologizados. Sean $\mathfrak J$, $V$, $W$ unas vecindades
abiertas de $0$ en $A$, $M$, $N$ respectivamente, que sean
$A$-módulos, y tales que $\mathfrak J.M\subset V$, $\mathfrak J.N\subset W$,
de modo que $M/V$ y $N/W$ pueden considerarse como
$(A/\mathfrak J)$-módulos. Cuando $\mathfrak J$, $V$, $W$ recorren los
sistemas de vecindades abiertas que verifican las condiciones anteriores, es
inmediato que los módulos
$(M/V)\otimes_{A/\mathfrak J}(N/W)$ forman un sistema proyectivo de módulos
\oldpage[0\textsubscript{I}]{76}
sobre el sistema proyectivo de anillos $A/\mathfrak J$; por paso al límite
proyectivo, se deduce, pues, un módulo sobre el separado completado $\widehat A$
de $A$, que se llama el \emph{producto tensorial completado} de $M$ y de
$N$ y que se denota por $(M\otimes_A N)^\wedge$. Si se observa que $M/V$
es canónicamente isomorfo a $\widehat M/\overline V$, donde $\widehat M$ es
el separado completado de $M$ y $\overline V$ la adherencia en $\widehat M$ de
la imagen de $V$, se ve que el producto tensorial completado
$(M\otimes_A N)^\wedge$ se identifica canónicamente con
$(\widehat M\otimes_{\widehat A}\widehat N)^\wedge$, que se denota también por
$\widehat M\widehat\otimes_{\widehat A}\widehat N$.
\end{env}

\begin{env}[7.7.2]
\label{0.7.7.2-fr}
\begin{sloppypar}
Con las notaciones anteriores, los productos tensoriales
$(M/V)\otimes_A(N/W)$ y
$(M/V)\otimes_{A/\mathfrak J}(N/W)$ se identifican canónicamente;
se identifican, por tanto, también con
$(M\otimes_A N)/(\Im(V\otimes_A N)+\Im(M\otimes_A W))$.
Se concluye que $(M\otimes_A N)^\wedge$ es el \emph{separado completado del
$A$-módulo $M\otimes_A N$, provisto de la topología para la cual los
submódulos}
\[
  \Im(V\otimes_A N)+\Im(M\otimes_A W)
\]
\emph{forman un sistema fundamental de vecindades de $0$} (recorriendo $V$ y $W$
el conjunto de los submódulos abiertos de $M$ y $N$
respectivamente); diremos, para abreviar, que esta topología es el
\emph{producto tensorial} de las topologías dadas sobre $M$ y $N$.
\end{sloppypar}
\end{env}

\begin{env}[7.7.3]
\label{0.7.7.3-fr}
Sean $M'$, $N'$ dos $A$-módulos linealmente topologizados,
$u:M\to M'$, $v:N\to N'$ dos homomorfismos continuos; es inmediato
que $u\otimes v$ es continuo para las topologías producto tensorial sobre
$M\otimes N$ y $M'\otimes N'$ respectivamente; por paso al separado
completado, se deduce un homomorfismo continuo
$(M\otimes N)^\wedge\to(M'\otimes N')^\wedge$, que designaremos
por $u\widehat\otimes v$; $(M\otimes_A N)^\wedge$ es, por tanto, un
\emph{bifuntor} en $M$ y $N$ en la categoría de los $A$-módulos
linealmente topologizados.
\end{env}

\begin{env}[7.7.4]
\label{0.7.7.4-fr}
Se define del mismo modo el producto tensorial completado de un número finito cualquiera
de $A$-módulos linealmente topologizados; es inmediato que este producto
posee las propiedades usuales de asociatividad y de conmutatividad.
\end{env}

\begin{env}[7.7.5]
\label{0.7.7.5-fr}
Si $B$, $C$ son dos $A$-álgebras topológicas linealmente topologizadas,
la topología producto tensorial sobre $B\otimes_A C$ tiene como sistema fundamental
de vecindades de $0$ los \emph{ideales}
$\Im(\mathfrak K\otimes_A C)+\Im(B\otimes_A\mathfrak L)$ del álgebra
$B\otimes_A C$, recorriendo $\mathfrak K$ (resp. $\mathfrak L$) el conjunto
de los ideales abiertos de $B$ (resp. $C$). Por consiguiente, $(B\otimes_A C)^\wedge$
está provisto de una estructura de \emph{$\widehat A$-álgebra topológica}, límite
proyectivo del sistema proyectivo de $(A/\mathfrak J)$-álgebras
$(B/\mathfrak K)\otimes_{A/\mathfrak J}(C/\mathfrak L)$ ($\mathfrak J$
ideal abierto de $A$ tal que $\mathfrak J.B\subset\mathfrak K$,
$\mathfrak J.C\subset\mathfrak L$; siempre existe uno). Se dice que
esta álgebra es el \emph{producto tensorial completado} de las álgebras $B$ y
$C$.
\end{env}

\begin{env}[7.7.6]
\label{0.7.7.6-fr}
Los homomorfismos de $A$-álgebras $b\to b\otimes 1$,
$c\to 1\otimes c$ de $B$ y $C$ en $B\otimes_A C$ son continuos
cuando se dota a esta última álgebra de la topología producto tensorial;
por composición con el homomorfismo canónico de $B\otimes_A C$ en su
separado completado, dan, pues, \emph{homomorfismos canónicos}
$\rho:B\to(B\otimes_A C)^\wedge$, $\sigma:C\to(B\otimes_A C)^\wedge$.
El álgebra $(B\otimes_A C)^\wedge$ y los homomorfismos $\rho$ y $\sigma$
poseen además la \emph{propiedad universal} siguiente: para toda
$A$-álgebra linealmente topologizada, separada y completa $D$, y todo
par de $A$-homomorfismos continuos $u:B\to D$, $v:C\to D$, existe un
$A$-homomorfismo continuo y uno solo
$w:(B\otimes_A C)^\wedge\to D$ tal que $u=w\circ\rho$ y
$v=w\circ\sigma$. En efecto, existe ya un $A$-homomorfismo único
$w_0:B\otimes_A C\to D$ tal que $u(b)=w_0(b\otimes 1)$ y
$v(c)=w_0(1\otimes c)$, y todo se reduce a probar que $w_0$ es \emph{continuo},
pues dará entonces un

\oldpage[0\textsubscript{I}]{77}

homomorfismo continuo $w$ por paso al separado completado. Ahora bien, si
$\mathfrak M$ es un ideal abierto de $D$, existen por hipótesis ideales
abiertos $\mathfrak K\subset B$, $\mathfrak L\subset C$ tales que
$u(\mathfrak K)\subset\mathfrak M$, $v(\mathfrak L)\subset\mathfrak M$;
la imagen por $w_0$ de
$\Im(\mathfrak K\otimes C)+\Im(B\otimes\mathfrak L)$ está también
contenida en $\mathfrak M$, de donde nuestra afirmación.
\end{env}

\begin{proposition}[7.7.7]
\label{0.7.7.7-fr}
Si $B$ y $C$ son dos $A$-álgebras preadmisibles,
$(B\otimes_A C)^\wedge$ es admisible, y si $\mathfrak K$ \textup{(resp.
$\mathfrak L$)} es un ideal de definición de $B$ \textup{(resp. $C$)}, la adherencia
en $(B\otimes_A C)^\wedge$ de la imagen canónica de
$\mathfrak H=\Im(\mathfrak K\otimes C)+\Im(B\otimes\mathfrak L)$ es
un ideal de definición.
\end{proposition}

Basta mostrar que $\mathfrak H^n$ tiende a $0$ para la topología
producto tensorial, lo que resulta inmediatamente de la inclusión
\[
  \mathfrak H^{2n}\subset
  \Im(\mathfrak K^n\otimes C)+\Im(B\otimes\mathfrak L^n)
\]

\begin{proposition}[7.7.8]
\label{0.7.7.8-fr}
Sean $A$ un anillo preádico, $\mathfrak J$ un ideal de definición de
$A$, $M$ un $A$-módulo \emph{de tipo finito}, provisto de la topología
$\mathfrak J$-preádica. Para toda $A$-álgebra topológica ádica y
noetheriana $B$, $B\otimes_A M$ se identifica con el producto tensorial completado
$(B\otimes_A M)^\wedge$.
\end{proposition}

Si $\mathfrak K$ es un ideal de definición de $B$, existe por hipótesis
un entero $m$ tal que $\mathfrak J^mB\subset\mathfrak K$, por tanto
$\Im(B\otimes_A\mathfrak J^{nm}M)=
\Im(\mathfrak J^{nm}B\otimes_A M)\subset
\Im(\mathfrak K^nB\otimes_A M)=\mathfrak K^n(B\otimes_A M)$; se
concluye que sobre $B\otimes_A M$, el producto tensorial de las topologías de $B$
y de $M$ es la topología $\mathfrak K$-preádica. Como $B\otimes_A M$
es un $B$-módulo de tipo finito, la proposición resulta inmediatamente de (7.3.6).

